\subsection{Instantané expérimental du 24 août 2026}
\label{sec:resultats-provisoires}

Les résultats présentés ici constituent un instantané de travail et non une
campagne finale. J'ai retenu quatre exécutions terminées et comparables au
regard des métadonnées disponibles, accessibles sur \texttt{nato-master}.
Elles utilisent toutes le modèle MiniMax-M2.7, le profil \texttt{full}, le
split \texttt{dev-public}, la révision logicielle \texttt{e7c4f6a} et la
politique d'évaluation \texttt{strict-v3} avec le contrat métrique
\texttt{strict-v3.3}. Les métadonnées indiquent que le modèle n'avait pas accès
à la vérité terrain. Chaque scénario ne comporte qu'une répétition ; les
valeurs ne permettent donc pas encore d'estimer la variance.

\begin{table}[H]
  \centering
  \small
  \caption{Résultats provisoires des quatre campagnes sélectionnées}
  \label{tab:resultats-provisoires}
  \begin{tabular}{lrrrrr}
    \toprule
    Scénario & Précision & Rappel & F1 & F1 vérifié & F1 qualité \\
    \midrule
    S1 & \(0{,}875\) & \(0{,}583\) & \(0{,}700\) & \(0{,}600\) & \(0{,}625\) \\
    S2 & \(0{,}833\) & \(0{,}769\) & \(0{,}800\) & \(0{,}640\) & \(0{,}670\) \\
    S5 & \(0{,}900\) & \(0{,}600\) & \(0{,}720\) & \(0{,}640\) & \(0{,}680\) \\
    S8 & \(0{,}533\) & \(0{,}571\) & \(0{,}552\) & \(0{,}414\) & \(0{,}448\) \\
    \midrule
    Moyenne macro & \(0{,}785\) & \(0{,}631\) & \(0{,}693\) & \(0{,}574\) & \(0{,}606\) \\
    \bottomrule
  \end{tabular}
\end{table}

Comme l'indique \cref{tab:resultats-provisoires}, la moyenne macro atteint un
F1 de \(0{,}693\), contre \(0{,}574\) lorsque seuls les constats vérifiés sont
crédités. Le F1 ajusté par la qualité s'établit à \(0{,}606\). L'écart entre
ces mesures suggère qu'une détection correcte au niveau du type et de la cible
ne s'accompagne pas toujours d'une preuve d'exécution suffisamment forte.

La complétion moyenne de la phase 4 atteint \(92{,}8\,\%\) et la couverture
d'exploitation \(82{,}4\,\%\). En revanche, la couverture macro des chemins
n'est que de \(52{,}1\,\%\), et celle des chemins entièrement vérifiés de
\(6{,}3\,\%\). S1 est le seul scénario de cet échantillon à obtenir une
couverture vérifiée non nulle, égale à \(25\,\%\). Dans cet échantillon, la
difficulté principale semble donc être la conservation d'une chaîne de preuves
complète sur plusieurs étapes, et non le seul lancement des vérifications
locales.

S2 obtient le meilleur F1 de l'échantillon avec \(0{,}800\). S8 est le cas le
plus difficile : huit vrais positifs, sept faux positifs et six faux négatifs
conduisent à un F1 de \(0{,}552\). Ce résultat motive les gardes consacrées à
la précision sur les scénarios étendus et à la réduction des hallucinations.

Les quatre exécutions totalisent environ 3,58 dollars, 11,0 millions de tokens,
1\,202 appels d'outils et 30 erreurs d'outils. Elles utilisent cependant un
modèle généraliste distant et non le \gls{hmoe} local. Cet instantané constitue
donc un test d'intégration du pipeline et du harness ; il ne démontre
pas encore l'efficacité des experts entraînés. Une campagne finale devra figer
la version du code, répéter chaque scénario et publier moyenne, dispersion et
intervalle de confiance pour les baselines et le \gls{hmoe}.

Les identifiants de runs utilisés sont
\texttt{2026-08-24\_092203} pour S1,
\texttt{2026-08-21\_144141} pour S2,
\texttt{2026-08-24\_083553} pour S5 et
\texttt{2026-08-21\_151305} pour S8.
